% =====================================================================
% Internship Report — RAGLab
% Document question-answering system for Arabic banking documents
%
% Compile with: pdflatex rapport_stage.tex (run twice).
% Arabic expressions are transliterated so this file compiles with
% pdfLaTeX without an Arabic font. Figures are drawn directly in LaTeX.
%
% Internship period: 3 August to 3 October 2026.
% Events from 4–5 October are isolated in a post-internship note.
% Each chapter, numbered or unnumbered, starts on a new page.
% =====================================================================
\documentclass[a4paper,12pt,openany]{report}

\usepackage[utf8]{inputenc}
\usepackage[T1]{fontenc}
\usepackage[english]{babel}
\usepackage{geometry}
\usepackage{setspace}
\usepackage{hyperref}
\usepackage{float}
\usepackage{array}
\usepackage{booktabs}
\usepackage{xcolor}
\usepackage{titlesec}
\usepackage{fancyhdr}

\geometry{margin=2.4cm,headheight=15pt}
\setstretch{1.3}
\setcounter{secnumdepth}{3}
\setcounter{tocdepth}{2}

% ============================================================
% Personal information to customize
% ============================================================
\newcommand{\StudentName}{[[INTERN_NAME]]}
\newcommand{\SchoolName}{Private Higher School of Engineering and Technology}
\newcommand{\HostOrganization}{Al Baraka Bank}
\newcommand{\SupervisorName}{Ahlem BENHADDOUD}
\newcommand{\InternshipPeriod}{3 August to 3 October 2026}
\newcommand{\AcademicYear}{2025 -- 2026}
\newcommand{\InternshipTopic}{Design and Evaluation of a Document Question-Answering
System (RAG) for Arabic Banking Documents}

\definecolor{genone}{HTML}{2F6FD0}
\definecolor{genpivot}{HTML}{B57A14}
\definecolor{gentwo}{HTML}{1F7A5C}
\definecolor{genlast}{HTML}{B3205A}

\hypersetup{
  hidelinks,
  pdftitle={Internship Report — RAGLab},
  pdfauthor={\StudentName},
  pdfsubject={RAG system for Arabic banking documents}
}

% With the report class, every chapter starts on a new page.
\titleformat{\chapter}[display]
  {\normalfont\bfseries\Huge}
  {\chaptername~\thechapter}
  {0.6em}
  {\Huge}
\titlespacing*{\chapter}{0pt}{0pt}{18pt}
\titleformat{\section}{\normalfont\bfseries\Large}{\thesection}{0.6em}{}
\titlespacing*{\section}{0pt}{2.2ex plus 0.5ex minus 0.2ex}{1.0ex}
\titleformat{\subsection}{\normalfont\bfseries\large}{\thesubsection}{0.6em}{}

\pagestyle{fancy}
\fancyhf{}
\fancyhead[L]{\small Internship Report — RAGLab}
\fancyhead[R]{\small\thepage}
\renewcommand{\headrulewidth}{0.3pt}
\fancypagestyle{plain}{%
  \fancyhf{}
  \fancyhead[L]{\small Internship Report — RAGLab}
  \fancyhead[R]{\small\thepage}
  \renewcommand{\headrulewidth}{0.3pt}
}

\begin{document}

% ====================== TITLE PAGE ======================
\begin{titlepage}
\centering
\vspace*{0.6cm}
{\large\textbf{Republic of Tunisia}}\\[0.2cm]
{\normalsize \SchoolName}\\[1.2cm]

{\Large\bfseries INTERNSHIP REPORT}\\[0.7cm]
{\large\bfseries \InternshipTopic}\\[1.3cm]

\begin{tabular}{>{\bfseries}p{4.2cm}p{9.4cm}}
Prepared by: & \StudentName \\[0.2cm]
Host organization: & \HostOrganization \\[0.2cm]
Internship period: & \InternshipPeriod \\[0.2cm]
Supervisor: & \SupervisorName \\[0.2cm]
Academic year: & \AcademicYear \\
\end{tabular}

\vfill
{\small Internship Report — \AcademicYear}
\end{titlepage}

% ====================== FRONT MATTER ======================
\pagenumbering{roman}

\chapter*{Dedication}
\addcontentsline{toc}{chapter}{Dedication}
I dedicate this work to my family, for their constant support throughout this journey, and to
everyone who encouraged me to turn an idea into a measured system.

\clearpage
\chapter*{Acknowledgements}
\addcontentsline{toc}{chapter}{Acknowledgements}
I would first like to thank \SupervisorName for her rigorous guidance, availability and scoping
decisions, which helped shape this work. I also thank the team at \HostOrganization for welcoming
me, providing access to the documents and sharing operational feedback that helped reveal defects
not visible in testing. Finally, I thank the teaching staff at \SchoolName for the education that
made this project possible.

\clearpage
\tableofcontents
\clearpage
\addcontentsline{toc}{chapter}{List of Figures}
\listoffigures
\clearpage
\pagenumbering{arabic}
\setcounter{page}{1}

% ====================== INTRODUCTION ======================
\chapter*{General Introduction}
\addcontentsline{toc}{chapter}{General Introduction}

Banks produce and consult a large volume of regulatory and procedural material, including laws,
central-bank circulars, internal guides and training documents. Finding reliable information in
these sources can be time-consuming. Giving this task to an unconstrained language model introduces
an additional risk: it may answer confidently even when the documents do not provide the necessary
evidence.

This internship, carried out from \textbf{3 August to 3 October 2026}, focused on the design,
implementation and evaluation of \textbf{RAGLab}, a document question-answering system based on
Retrieval-Augmented Generation (RAG) and applied to a real Arabic banking corpus. The project had
three main objectives:

\begin{enumerate}
  \item \textbf{Traceability:} every claim must be supported by a verbatim excerpt from the cited
document, and every stated number must appear in that evidence.
  \item \textbf{Honest refusal:} when the available evidence is insufficient, the system must refuse
to answer and explain what is missing.
  \item \textbf{Measured evaluation:} architectural choices must be compared using an explicit,
reproducible protocol rather than selected on the basis of impressions alone.
\end{enumerate}

The report is organized into four numbered chapters. Chapter 1 presents the project context and
scope; Chapter 2 explains the RAG principles and technical choices; Chapter 3 describes the
architecture and its evolution; Chapter 4 analyses the results, limitations and operating
conditions. A conclusion presents the findings and future work. Changes dated \textbf{4 and
5 October 2026}, after the internship ended, are isolated in a separate note so they are not
mistaken for work completed during the internship.

The reported figures are linked to measurement campaigns or to sources listed in the appendices.
Retrieval results expressed as \texttt{hit@k} are reported as percentages.

% ====================== CHAPTER 1 ======================
\chapter{Project Context and Scope}

\section{Host Organization and Work Environment}
The internship took place at \textbf{Al Baraka Bank}, an Islamic bank in Tunisia. The banking
activities and products relevant to this project fall, among other things, within the framework of
\textbf{Law No. 2016-48 of 11 July 2016 on banks and financial institutions}, which is one of the
documents in the corpus.

The project was carried out in an individual development environment using Python 3.11, a Git
repository, a locally containerized service and continuous integration (CI), which served both as
a quality gate and as a benchmark for real models. Two hardware and network constraints shaped the
method: model-provider APIs were not accessible from the development workstation, and no GPU was
available. Live measurements were therefore run in CI, and models were accessed through APIs.

\section{Corpus and Business Problem}
The evaluation corpus contains four real documents, all in Arabic:

\begin{table}[htbp]
\centering
\small
\begin{tabular}{p{4.6cm}p{2.9cm}p{6.4cm}}
\toprule
\textbf{Document} & \textbf{Type} & \textbf{Role in the Corpus} \\
\midrule
Law No. 2016-48 (11 July 2016) & Law & Primary legal reference (198 articles, 11 titles) \\
BCT Circular No. 2019-08 & Circular & Operational guidance (20 chapters) \\
Internal Guide to Islamic Banking Operations & Internal guide & Internal compliance policy and
procedures \\
Introduction to Islamic Finance (\emph{Madkhal}) & Training material & History, definitions and
orders of magnitude \\
\bottomrule
\end{tabular}
\end{table}

Three practical challenges motivated the project:
\begin{itemize}
  \item \textbf{A mismatch between the user's vocabulary and the document's vocabulary.} A user may
ask, ``Can I open a business?'', while the corpus uses technical terms such as \emph{al-sigha},
\emph{mahall al-aqd} or \emph{tamwil}.
  \item \textbf{Asymmetric multilingualism.} The corpus is in Arabic, while questions may be asked
in Arabic, French or English. A French question should be able to retrieve an Arabic passage
without automatically translating either the document or the query.
  \item \textbf{The cost of an error is high.} An incorrect answer about a rate, deadline or penalty
may have legal consequences; linguistic plausibility alone is not enough.
\end{itemize}

\section{Internship Objectives}
\begin{enumerate}
  \item Build a complete pipeline: loading, chunking, embedding, indexing, retrieval, grounded
generation and verification.
  \item Compare accessible models and document their performance on the real corpus.
  \item Make explicit refusal a first-class behaviour, supported by a taxonomy of reasons.
  \item Deliver the system as a service with a stable, versioned and tested HTTP contract.
  \item Instrument quality with offline tests, frozen question sets and an audit log.
\end{enumerate}

\section{Project Constraints}
\begin{table}[htbp]
\centering
\small
\begin{tabular}{p{6.0cm}p{8.2cm}}
\toprule
\textbf{Constraint} & \textbf{Methodological Consequence} \\
\midrule
No network access to model APIs from the workstation & Measurements on real models run in CI and
are triggered explicitly to control cost \\
No GPU & Use of API-hosted models, pinned and versioned \\
Real corpus only & No synthetic documents in the final corpus; test sets are frozen and versioned \\
Service under integration & HTTP endpoints are frozen; changes are additive and the contract is
versioned \\
Traceability requirement & Each result must be linked to a commit, log or execution identifier \\
\bottomrule
\end{tabular}
\end{table}

\section{Working Method}
After the methodological pivot described in Chapter 3, the project followed a protocol called
\textbf{“step and plan”}. Before each step, a plan specifies the objective and scope; inputs;
actions; outputs; automated checks; how the sponsor can test the result; the acceptance criterion;
and the rollback point.

Two rules complement this protocol. The local gate (compilation, offline checks, unit tests and
inspection) must be green \textbf{before} publication; CI is then monitored until it is green as
well. The question sets remain frozen: a campaign that modifies them is no longer directly
comparable with earlier campaigns.

Implementation was carried out with the assistance of a coding agent under a \textbf{written work
contract} recording constraints, verified facts and previously encountered pitfalls. Every step
remained subject to human validation: the agent proposes, measures and documents, while activation
decisions remain with the project sponsor. This organization made it easier to resume work after
interruptions and reduced repeated diagnosis of the same defects.

% ====================== CHAPTER 2 ======================
\chapter{Background and Technical Choices}

\section{RAG Principles and Main Risks}
A RAG system combines an information-retrieval engine, which selects passages from a document
collection, and a language model, which drafts an answer from those passages. This architecture
does not eliminate risk; it shifts risk to the quality of chunking, retrieval, evidence assessment
and answer validation.

\begin{table}[htbp]
\centering
\small
\begin{tabular}{p{3.5cm}p{5.0cm}p{5.7cm}}
\toprule
\textbf{Risk} & \textbf{Symptom} & \textbf{Mitigation} \\
\midrule
Arbitrary chunking & A definition is split and becomes unusable & Structural chunking (Section 3.4.1) \\
Lexical mismatch & The query and document use different forms & Arabic normalization, bridges
between writing systems and bounded interrogation \\
Similarity false positives & Irrelevant passages are returned without warning & Deterministic
sufficiency state and explicit refusal \\
Hallucination & A plausible answer introduces a number absent from the source & Citation and
number checks (Section 3.4.4) \\
Loss of traceability & The passage used by the model cannot be verified & Verbatim citations,
fragment identifiers and audit log \\
Silent regressions & A local improvement degrades another case & Frozen test sets and paired
measurements \\
\bottomrule
\end{tabular}
\end{table}

\section{Chunking: Fixed Windows and Document Structure}
\textbf{Fixed-window chunking} splits text every $N$ tokens, with overlap. It is simple but ignores
document structure. \textbf{Structural chunking} respects boundaries such as headings, articles,
paragraphs and tables; it requires a document-reconstruction step beforehand.

Both modes were implemented and compared. The historical \texttt{size} mode uses 220 tokens with
an overlap of 40. The \texttt{restructure} mode, which became the default, follows three steps:
\begin{enumerate}
  \item \textbf{Semantic normalization:} clean gazette headers, remove repeated headers, extract
section markers (\emph{al-fasl}, \emph{al-unwan}, \emph{al-bab}) and reconstruct the visual order
of Arabic lines extracted from PDFs.
  \item \textbf{Breadcrumb enrichment:} add a context line of the form \texttt{h1 > h2 > h3} above
each subsection heading and table.
  \item \textbf{Recursive structural chunking:} split hierarchically on headings, blank lines, line
breaks and spaces, with a target budget of 220/40 tokens.
\end{enumerate}
Linguistic repair of extracted text complements this process: some Arabic legal PDFs are extracted
in visual order, which can corrupt numbers in particular.

\section{Vector Representations}
Multilingual embedding models project texts into a space where vector proximity represents semantic
proximity. For Arabic, rich morphology and the coexistence of different writing systems make it
harder to match a question with a document passage.

Several providers were compared (Google, Jina, HuggingFace and NVIDIA). The selected embedding
model is \texttt{nvidia/nemotron-3-embed-1b}, with a native dimension of 2048. Each index stores a
fingerprint of the model and chunking configuration that produced it; a search using an
incompatible fingerprint is rejected before the model is called.

\section{Dense, Lexical and Hybrid Retrieval}
Three strategies were compared:
\begin{itemize}
  \item \textbf{Dense retrieval:} cosine similarity between vectors.
  \item \textbf{RRF:} reciprocal-rank fusion of dense retrieval and the BM25 lexical index
(constant 60; $k_1=1.5$; $b=0.75$).
  \item \textbf{Blend:} score fusion, $\lambda \cdot \text{cosine} + (1-\lambda)\cdot
\text{normalized BM25}$, with $\lambda=0.7$.
\end{itemize}
The results in Chapter 4 showed no net gain from lexical fusion on this corpus. Dense retrieval
was therefore retained as the default; the other modes remain explicitly available.

\section{Grounded Answers and Evidence Sufficiency}
The model must return a \texttt{grounded-v1} JSON object rather than free-form text. Each claim must
be associated with one or more verbatim citations from the supplied passages. Three deterministic
checks are applied:
\begin{enumerate}
  \item \textbf{Citation membership:} after normalization (Arabic, case and whitespace), each
citation must appear in the referenced passage.
  \item \textbf{Numeric check:} every number in a claim must appear in its supporting citations,
including thousands separators and Eastern Arabic numerals.
  \item \textbf{Policy check:} a cited source must have been retrieved, belong to the authorized
documents and, when validity metadata is available, be in force.
\end{enumerate}

Before generation, a deterministic \textbf{sufficiency} layer compares the evidence required by the
query intent with the content actually retrieved. It produces four states: sufficient,
insufficient, contradictory or indeterminate. This state, not a similarity score, determines
whether the system answers, refuses or attempts one bounded reinterpretation.

\section{Evaluation Protocol}
\textbf{Retrievability} is measured with \texttt{hit@1}, \texttt{hit@3} and \texttt{hit@5}: does
the expected document and passage appear among the top $k$ results? Final behaviour is also
examined, including grounded answers, reasoned refusals and citation compliance.

Two frozen test sets serve as references: a 50-case set (10 verbatim, 20 paraphrases,
15 cross-document and 5 out-of-scope cases) and a 10-question Arabic set covering five target
categories: colloquial language, synonyms, implicit requests, compound requests and ambiguity.

Comparisons are measured within the same run; a tolerance of $\pm 2$ percentage points is allowed
between runs because hosted models are not bit-reproducible. Frozen test sets are never modified to
improve a result.

\section{Selected Models}
Development checks, blind tests and prompt-injection scenarios led to the following model pair:

\begin{table}[htbp]
\centering
\small
\begin{tabular}{p{2.6cm}p{5.3cm}p{6.3cm}}
\toprule
\textbf{Role} & \textbf{Provider / Model} & \textbf{Rationale} \\
\midrule
Embeddings & NVIDIA \texttt{nemotron-3-embed-1b} (2048 dimensions) & Measured quality on the
Arabic corpus, high native dimension and caching support \\
Answer generation & xKiro \texttt{qwen/qwen3.8-max:free} & Highest rate of citations compliant
with the \texttt{grounded-v1} contract; live cost checks \\
\bottomrule
\end{tabular}
\end{table}

No fallback to another model is allowed: a different model causes an explicit error rather than a
silent degradation. Automatic query translation, which was measured and documented as a historical
approach, was removed.

\section{Technical Stack}
\begin{table}[htbp]
\centering
\small
\begin{tabular}{p{3.2cm}p{5.6cm}p{5.4cm}}
\toprule
\textbf{Component} & \textbf{Choice} & \textbf{Alternative Rejected} \\
\midrule
Orchestration & Explicit Python code & LangChain / LlamaIndex: opaque behaviour and dependencies \\
Vector database & Local persistent ChromaDB, cosine distance & Cloud database: cost and banking data \\
HTTP calls & Standard-library \texttt{urllib} & Provider SDKs \\
Service & FastAPI + uvicorn & — \\
Chunking & \texttt{tiktoken} (\texttt{cl100k\_base}) with a fallback estimate & — \\
PDF / DOCX & \texttt{pypdf} and direct OOXML reading & Office SDKs \\
Tests / CI & \texttt{unittest} + GitHub Actions & — \\
\bottomrule
\end{tabular}
\end{table}

The repository uses no orchestration framework, external vector database or provider SDK. This
choice simplifies auditing and deployment on an internal host.

% ====================== CHAPTER 3 ======================
\chapter{Architecture and Implementation}

\section{Overview: Two Generations and a Pivot}
The architecture evolved during the project. The first generation primarily addressed the question
of how to compare retrieval methods. After a methodological pivot, the second generation added
explicit layers for knowledge, understanding, control and auditing.

\begin{figure}[H]
\centering
\begin{tabular}{c}
\colorbox{genone!10}{\parbox{0.86\textwidth}{\centering
  \textbf{Generation 1 — Information Retrieval Laboratory}\\
  Test corpus, multiple providers, query variants, chunking experiments and retrieval fusion}}
  \\[0.45em]
$\Downarrow$ \textit{methodological pivot: 24--28 September 2026}\\[0.45em]
\colorbox{gentwo!10}{\parbox{0.86\textwidth}{\centering
  \textbf{Generation 2 — Governed Four-Layer Architecture}\\[0.2em]
  1. Corpus and structure \quad 2. Knowledge \quad 3. Understanding \quad 4. Delivery and audit}}
\end{tabular}
\caption{Evolution of the RAGLab architecture during the internship.}
\label{fig:architecture}
\end{figure}

\begin{figure}[H]
\centering
\begin{tabular}{p{3.6cm}p{10.2cm}}
\toprule
\textbf{Period} & \textbf{Milestone} \\
\midrule
4--23 September & Generation 1: flat, multi-provider pipeline; translation experiments and A/B
comparisons \\
24--28 September & Pivot: target-state report, gap analysis, refocusing on the real corpus and
adoption of the “step and plan” protocol \\
24 September--3 October & Generation 2: four-layer architecture, knowledge units and sufficiency
controls \\
4--5 October & Post-internship changes, described separately in the complementary note \\
\bottomrule
\end{tabular}
\caption{Documented milestones. The 4--5 October milestones fall outside the internship period.}
\label{fig:chronology}
\end{figure}

The engine shared by both generations remains deliberately simple: NVIDIA embeddings, a local
cosine index, answer generation by the pinned model and citation checks. Most of the evolution
concerns the corpus, chunking, and the knowledge, understanding and traceability layers added
around that engine.

\section{Generation 1: Information Retrieval Laboratory}
The first generation was a comparison-oriented pipeline: test corpus; \texttt{size} and then
\texttt{restructure} chunking; multi-provider embeddings and caching; ChromaDB indexing;
\texttt{vector}/\texttt{rrf}/\texttt{blend} retrieval; automatically translated query variants;
multi-provider answers; and citation checks.

This generation answered the question, “Which chunking method, model and retrieval fusion perform
best?” It did not explicitly determine whether the evidence was sufficient, whether the system
should refuse, or how decisions should be traced. These limitations motivated the pivot.

\section{Methodological Pivot (24--28 September 2026)}
The pivot had four steps:
\begin{enumerate}
  \item \textbf{Reverse specification:} static code inspection, with each claim cited (file,
symbol and line) and labelled as observed, inferred or unknown.
  \item \textbf{Gap analysis:} comparison of the observed system with the target-state report.
Citation checks, refusal categories and rejection of stale indexes already existed; intent,
evidence planning, sufficiency, document governance and audit logging needed further work.
  \item \textbf{Refocusing on the real corpus:} the web-compiled campaign was retired and the four
real banking documents became the working corpus.
  \item \textbf{Corpus repair:} 215 corrections to the law (29 batches), 35 to the circular
(including 28 corrupted numbers), 23 to the guide and 21 to the training material. Each log was
checked automatically and each document approved by the project sponsor. The local evaluation
record then contained 339 chunks.
\end{enumerate}

\section{Generation 2: Layered Architecture}

\subsection{Layer 1 — Corpus and Structure}
\begin{table}[htbp]
\centering
\small
\begin{tabular}{p{3.0cm}p{11.5cm}}
\toprule
\textbf{Module} & \textbf{Role} \\
\midrule
\texttt{loader.py} & PDF/DOCX/MD loading; Arabic normalization (NFKC, presentation forms,
alif normalization, tatweel and diacritic removal); language detection \\
\texttt{restructure.py} & Structural chunking; visual-order repair through zone reconstruction;
adoption of corrected texts \\
\texttt{chunker.py} & Historical \texttt{size} chunking (220/40); chunking fingerprint recording
mode, size, overlap and tokenizer \\
\texttt{docstore.py} & Application-pushed documents versioned by SHA-256 fingerprint; index marked
stale until re-indexing \\
\texttt{store.py} & Local cosine ChromaDB, BM25 lexical index, RRF/blend fusion, consistency checks
and source-level deletion \\
\bottomrule
\end{tabular}
\end{table}

Two safeguards structure this layer. The index is \textbf{self-describing}: each chunk records the
fingerprint of its chunking configuration and embedding space, and searches against an
incompatible fingerprint are rejected. An \textbf{indexing lock} prevents inconsistent reads
during ingestion and returns an explicit conflict instead.

\subsection{Layer 2 — Knowledge}
Built on Law 2016-48, this layer is deterministic and declarative: its data files are versioned,
automatically validated and reviewable line by line.

\begin{table}[htbp]
\centering
\small
\begin{tabular}{p{3.4cm}p{7.4cm}p{3.7cm}}
\toprule
\textbf{Element} & \textbf{Content} & \textbf{Check} \\
\midrule
Knowledge units & 198 units corresponding to the articles, with verbatim text, stable identifiers
such as \texttt{loi-2016-48:artNNN}, hierarchical paths and governed functional types & 198/198
coverage, literal text and stable identifiers \\
Relations & 77 internal references and 1 anchor relation (circular to Article 11) & Verbatim evidence
required; targets must exist \\
Structured numbers & 53 records (14 rates, 28 deadlines, 5 caps, 6 penalties) and a correction table;
served by \texttt{GET /numbers} & Each value must appear literally in its unit \\
Governance & 6 axes: authority, validity, scope, audience, domain and source & Registry-to-corpus
consistency \\
Unit index & 198 descriptive rows in an isolated collection & Law-question measurement:
100/100/100 vs. 67/92/92\% \\
\bottomrule
\end{tabular}
\end{table}

The addition of \texttt{GET /numbers} is an additive extension: read-only, filterable, usable even
with an empty index, and it does not change existing endpoints.

\subsection{Layer 3 — Understanding the Request}
The design rule is: \textbf{deterministic first, model second, and never twice}.

\begin{table}[htbp]
\centering
\small
\begin{tabular}{p{4.3cm}p{7.7cm}p{2.5cm}}
\toprule
\textbf{Module} & \textbf{Function} & \textbf{Status at Internship End} \\
\midrule
\texttt{intent.py} & Declarative intent classification; detection of compound requests and
influential ambiguity & Active \\
\texttt{evidence\_plan.py} & Derives evidence requirements from intent, never the answer & Active \\
\texttt{sufficiency.py} & Compares the plan with retrieved passages; identifies sufficiency,
insufficiency and contradictions; bounds guided retrieval turns & Active \\
\texttt{decompose.py} & Breaks a question into micro-questions with one evidence requirement each &
Built, not wired in \\
\texttt{interrogate.py} & If nothing is covered, converts a practical request into a technical
question in one bounded call & Active \\
\texttt{relational\_expansion.py} & Expands retrieval along relations, appended to the results &
Disabled (measured as neutral) \\
\texttt{micro\_retrieval.py} & Micro-question retrieval with result fusion & Disabled (measured as
unfavourable) \\
\bottomrule
\end{tabular}
\end{table}

The topic map associated with interrogation was reconstructed on 4--5 October. Since this change
occurred after the internship, it is described separately in the complementary note.

Interrogation has three safeguards: \textbf{fail closed} (malformed output, invented topics or
requirements, or an unchanged reformulation lead to refusal); \textbf{full disclosure} (the
\texttt{understood\_as}, \texttt{original\_question} and \texttt{interrogation} fields); and
\textbf{visible cost} (the call duration is exposed separately).

\subsection{Layer 4 — Delivery and Audit}
\begin{itemize}
  \item \textbf{Generation:} the \texttt{grounded-v1} prompt treats the question and passages as
untrusted data and requires strict JSON output (claims and evidence).
  \item \textbf{Citation checks:} citation membership, numbers and policy are validated. Law
passages carry stable knowledge-unit identifiers.
  \item \textbf{Sufficiency enforcement:} if no evidence requirement is covered, the service refuses
before generation; partial coverage allows one bounded regeneration.
  \item \textbf{Scrubbing:} recognizable personal data (email, phone, Tunisian RIB/IBAN, national ID)
are replaced with labels after validation.
  \item \textbf{Audit log:} each request produces a JSONL record (identifier, timestamp, scrubbed
question, status, sufficiency state, model, counters and latency), available through
\texttt{GET /audit} with bounded retention.
\end{itemize}

\subsection{Request Lifecycle}
\begin{verbatim}
Question (ar / fr / en)
  -> authentication, greeting shortcut, indexing lock
  -> intent -> evidence plan
  -> dense retrieval on the ORIGINAL question (never translated), top 5
  -> sufficiency check on retrieved passages
       |- sufficient -> direct generation (zero additional calls)
       |- partial    -> one bounded regeneration for covered points
       '- none       -> ONE bounded interrogation call -> new retrieval -> same check
  -> context construction (3000-token budget; truncated passage is discarded)
  -> grounded JSON generation -> citation checks -> scrubbing -> audit log
  -> cited answer, or reasoned refusal (both are valid responses)
\end{verbatim}

\subsection{Feature Gates and Defaults}
Experimental components are controlled by explicit feature gates. During the period described,
the active gates covered sufficiency fields, sufficiency enforcement, interrogation, concept
vocabulary and bridges between writing systems. Relational expansion and micro-question retrieval
remained disabled based on measurements. Each disabled gate restores the prior behaviour, as
verified by a dedicated test.

\section{Service and Interfaces}
\begin{table}[htbp]
\centering
\small
\begin{tabular}{p{3.4cm}p{11.1cm}}
\toprule
\textbf{Aspect} & \textbf{Implementation} \\
\midrule
Endpoints & 27 routes: health, profiles and models, keys, inspection, chunks, indexing, retrieval,
answers, evaluation, documents, diagnostics, audit and numbers \\
Version & \texttt{SERVICE\_VERSION = 1.4.0} at the end of the internship; version 1.5.0 is
post-internship \\
Contract & Human-readable contract document and generated OpenAPI schema \\
Authentication & Shared \texttt{X-Service-Token}, constant-time comparison; service intended to run
behind the application gateway \\
Errors & Single envelope \texttt{\{"detail": \{"reason": ...\}\}}, stable reasons and filtered
sensitive values \\
Evolution & Existing endpoints frozen; changes are additive \\
Integration & CLI, local console and HTTP console for integration testing \\
\bottomrule
\end{tabular}
\end{table}

The service is delivered with a \texttt{Dockerfile} and \texttt{docker-compose.yml}, including
volumes for the index, embedding cache and documents pushed by the application. The integration
guide covers the HTTP contract, reference queries, interface mock-up and deployment.

\section{Quality Engineering}
\begin{table}[htbp]
\centering
\small
\begin{tabular}{p{3.4cm}p{11.1cm}}
\toprule
\textbf{Mechanism} & \textbf{Content} \\
\midrule
Offline tests & Unit tests, deterministic checks, compilation, inspection and dependency checks \\
Continuous integration & Offline gate on every push; manually triggered real-model campaigns; heavy
harness, catalogs, retrieval judge and answer A/B tests \\
Quality gates & Index fingerprints, refusal during indexing, no fallback, refusal before generation
and fail-closed interrogation \\
Reproducibility & Fingerprint-addressed caches; archived runs; results linked to a commit and run
identifier \\
Documentation & Contract, query guide, interface mock-up, deployment guide and work contract for
follow-up sessions \\
\bottomrule
\end{tabular}
\end{table}

% ====================== CHAPTER 4 ======================
\chapter{Evaluation, Results and Limitations}

\section{Corpus and Question Sets}
The evaluation corpus contains four real documents that were repaired and approved. The local
record reports \textbf{339 chunks} for the structural arm. CI produces 858 with the full tokenizer;
this difference is related to the tokenization environment. These volumes should not be compared as
if they came from the same counting setup.

\begin{table}[htbp]
\centering
\small
\begin{tabular}{p{3.6cm}p{3.0cm}p{7.9cm}}
\toprule
\textbf{Set} & \textbf{Size} & \textbf{Composition} \\
\midrule
\texttt{questions\_50.json} & 50 cases (45 evaluable, 5 out of scope) & 10 verbatim,
20 paraphrases, 15 cross-document (6 fr, 9 en); 32 ar, 8 fr and 10 en \\
\texttt{questions\_targets.json} & 10 Arabic cases & 2 per category: colloquial language,
synonyms, implicit, compound and ambiguous \\
\bottomrule
\end{tabular}
\end{table}

Every expected substring was checked automatically: it had to appear literally in the expected
document and be distinctive within the corpus. The main test set contains 12 law cases, 10 circular
cases, 11 guide cases and 12 training-material cases.

\section{Measurement Protocol}
\begin{itemize}
  \item \textbf{Paired comparison:} compared arms are measured in the same run.
  \item \textbf{Tolerance:} $\pm 2$ percentage points for comparisons across runs.
  \item \textbf{Two stages:} deterministic, free BM25-only retrieval serves as a quick filter;
embeddings and real models decide the live-stage results.
  \item \textbf{Final behaviour:} response status (grounded answer or refusal) and citation
compliance are examined in addition to \texttt{hit@k}.
\end{itemize}

\section{Chunking: Fixed Windows vs. Structural Chunking}
Offline measurement (BM25 only, $k=20$) on the 50-case set:

\begin{table}[htbp]
\centering
\begin{tabular}{lccc}
\toprule
\textbf{Arm} & \textbf{hit@1 (\%)} & \textbf{hit@3 (\%)} & \textbf{hit@5 (\%)} \\
\midrule
\texttt{size} 220/40 & 42 & 60 & 62 \\
\texttt{restructure} (repaired, approved corpus) & \textbf{62} & \textbf{82} & \textbf{87} \\
\bottomrule
\end{tabular}
\end{table}

By language, results were: Arabic 43/57/60 then 67/87/93; English 44/89/89 then 67/100/100;
French 33/33/33 in both arms. The lexical ceiling in French motivated vector-based evaluation.

\textbf{Live} measurement (embeddings and real models, same question set):
\begin{table}[htbp]
\centering
\begin{tabular}{lccc}
\toprule
\textbf{Arm} & \textbf{hit@1 (\%)} & \textbf{hit@3 (\%)} & \textbf{hit@5 (\%)} \\
\midrule
\texttt{size} & 51 & 64 & 71 \\
\texttt{restructure} & \textbf{71} & \textbf{80} & \textbf{87} \\
\bottomrule
\end{tabular}
\end{table}

The gain is \textbf{20 percentage points in hit@1}. Some cases missed by dense retrieval were
retrieved by the lexical arm, motivating a deterministic reranker that was evaluated separately.

\section{Retrieval Modes: Fusion Does Not Pay Off}
\begin{table}[htbp]
\centering
\small
\begin{tabular}{lcccl}
\toprule
\textbf{Mode} & \textbf{hit@1 (\%)} & \textbf{hit@3 (\%)} & \textbf{hit@5 (\%)} &
\textbf{Decision} \\
\midrule
\texttt{vector} & \textbf{71} & 80 & \textbf{87} & Retained as default \\
\texttt{rrf} (BM25 fusion) & 71 & 80 & 82 & Tied at rank 1, lower recall \\
\texttt{blend} ($\lambda=0.7$) & 62 & 80 & 84 & Lower at rank 1 \\
\bottomrule
\end{tabular}
\end{table}

In French, dense retrieval scores 50/67/67, compared with 33/33/50 for RRF and 33/33/67 for the
blend. The decision rule is to change the default only if a candidate wins without regressions on
verbatim and out-of-scope cases. Dense retrieval was therefore retained; the other modes remain
explicit, revisable options.

\section{Answer-Model Independence}
\begin{table}[htbp]
\centering
\small
\begin{tabular}{lcccc}
\toprule
\textbf{Model} & \textbf{Answers} & \textbf{Refusals} & \textbf{Errors} &
\textbf{Gate refusals} \\
\midrule
\texttt{qwen/qwen3.8-max:free} (selected) & 37 & 13 & 0 & 7 \\
\texttt{moonshotai/kimi-k3} & 27 & 14 & 9 & 7 \\
\bottomrule
\end{tabular}
\end{table}

Status agreement between the two models is \textbf{0.720}. When both models answer, their cited
sources are identical (Jaccard similarity = 1.000). With the context fixed, model disagreement is
therefore mainly about whether to answer or refuse, rather than which passages to cite. This
shifts the remaining work toward evidence completeness.

\section{Knowledge-Layer Results}
On the 12 law questions, the unit index was compared with structural chunking:

\begin{table}[htbp]
\centering
\begin{tabular}{lccc}
\toprule
\textbf{Arm} & \textbf{hit@1 (\%)} & \textbf{hit@3 (\%)} & \textbf{hit@5 (\%)} \\
\midrule
Reference (structural chunking) & 67 & 92 & 92 \\
Unit index & \textbf{100} & \textbf{100} & \textbf{100} \\
\bottomrule
\end{tabular}
\end{table}

The ten target Arabic questions produced the following results:
\begin{table}[htbp]
\centering
\small
\begin{tabular}{lccc}
\toprule
\textbf{Category} & \textbf{hit@1 (\%)} & \textbf{hit@3 (\%)} & \textbf{hit@5 (\%)} \\
\midrule
Compound request & 100 & 100 & 100 \\
Implicit and ambiguous & 50 & 100 & 100 \\
Synonyms & 50 & 50 & 50 \\
Colloquial language & 0 & 0 & 0 \\
\bottomrule
\end{tabular}
\end{table}

Activating a 16-entry institutional lexicon brought no gain to dense retrieval, although it
improved the lexical arm. The lexicon was therefore left disabled for the dense default.

\section{Robustness}
Six adversarial scenarios are checked automatically by the citation validator: fabricated
citation, citation borrowed from another source, calculated number, unknown citation identifier,
citation marker embedded in a claim, and instruction injection in either the question or a
passage. All six are rejected by the deterministic validator.

An important limitation remains: the validator guarantees that words and numbers come from the
source, but not that a claim logically follows from it. Other operational defects were corrected
during the project, with regression tests: network errors leaking as HTTP 500 (replaced by an
explicit 502), locking during ingestion, malicious filenames and container configuration in which
a hard-coded token silently overrode environment variables.

\section{Limitations}
\begin{table}[htbp]
\centering
\small
\begin{tabular}{p{4.7cm}p{9.8cm}}
\toprule
\textbf{Limitation} & \textbf{Impact} \\
\midrule
Citation checking $\neq$ entailment & A correctly sourced sentence can still be an abusive
interpretation \\
No similarity threshold & Five passages are passed on even if irrelevant; safety relies on
sufficiency checks and refusal \\
Single process & One profile, one index and one indexing operation at a time; no horizontal scaling \\
Updates require re-indexing & Replaced documents may continue to serve old chunks until
re-indexing \\
Tokenizer & Without the \texttt{cl100k} vocabulary, token counts are estimates; CI is authoritative \\
Heuristic language detection & Misclassification remains possible \\
No conversation memory or streaming & Deliberate choices in the current configuration \\
Small, correlated test sets & 50 + 10 cases provide neither a security guarantee nor a legal
assurance \\
\bottomrule
\end{tabular}
\end{table}

The system is \textbf{not ready for production banking use}. It may be suitable for a supervised
pilot using approved, non-confidential documents, subject to the required validations.

\section{Cost and Operations}
Fingerprint-based caching and free models keep offline campaign costs very low. Live campaigns are
manually triggered and bounded. Each campaign publishes its figures alongside its run. Container
packaging, persistent volumes, token-aware health checks and a deployment guide were prepared for
the platform team.

% ====================== CONCLUSION ======================
\chapter*{Conclusion and Future Work}
\addcontentsline{toc}{chapter}{Conclusion and Future Work}

\section*{Outcome}
The internship produced a measured and documented document question-answering system designed for
a real Arabic banking corpus:
\begin{itemize}
  \item a complete pipeline, from document reading and repair through retrieval, grounded
generation, citation checks and reasoned refusal;
  \item a layered architecture (structure, knowledge, understanding, delivery and audit) resulting
from an intentional change in method;
  \item quantified results, including a 20-percentage-point gain in hit@1 from structural chunking
and a 100/100/100 score for the unit index on law questions;
  \item a documented, versioned service with an HTTP contract, audit log, personal-data scrubbing
and containerized configuration.
\end{itemize}

\section*{Lessons Learned}
\begin{enumerate}
  \item \textbf{Measurement replaces opinion.} Chunking, retrieval and feature-activation choices
were decided through paired comparisons.
  \item \textbf{Governance is part of the architecture.} Plans, feature gates and frozen test sets
reduced the risk of silent regressions.
  \item \textbf{The real service remains essential.} Some defects appear only in operation, so
testing must be complemented by checks against the deployed service.
\end{enumerate}

\section*{Future Work}
\begin{enumerate}
  \item Finalize identity and role decisions before use beyond an internal pilot.
  \item Add topic metadata to chunks (\texttt{path}, \texttt{level}, \texttt{parent\_id}) in an
isolated collection before any production re-indexing.
  \item Measure faceted retrieval and any fusion with dense retrieval.
  \item Re-run all campaigns on frozen test sets and document regressions.
  \item Prepare for industrialization: multiple users, incremental updates, observability and
continuous adversarial evaluation.
\end{enumerate}

% ====================== POST-INTERNSHIP NOTE ======================
\chapter*{Post-Internship Update — 4--5 October 2026}
\addcontentsline{toc}{chapter}{Post-Internship Update — 4--5 October 2026}

\emph{Scope.} The items below are dated \textbf{4 and 5 October 2026}, after the internship end
date of 3 October. They are retained to document the project's subsequent evolution, but are not
counted as work completed during the internship.

\section*{Phase 9: Topic Map and Operational Fixes}
The topic map, previously derived only from section titles, was rebuilt from source-backed units:
262 nodes, including 198 law units and 20 circular chapters, each with an identifier, path and
excerpt. Interrogation now uses exact-match identifiers. The service moved to version 1.5.0 without
changing the response format.

Post-internship measurements produced the following results:
\begin{table}[htbp]
\centering
\small
\begin{tabular}{p{6.2cm}p{8.3cm}}
\toprule
\textbf{Measurement} & \textbf{Result} \\
\midrule
Retrieval only, 7 evidence-bearing cases & \texttt{hit@1/3/5}: 5/7 for dense retrieval \\
Topic-map comparison, 12 live calls & 12/12 parsed without error; target topics reached rose
from 4/5 to \textbf{5/5} \\
End-to-end probe, 6 unchanged questions & \textbf{3 grounded answers and 3 refusals}; all
generated answers validated; no provider errors \\
Number of retrieved chunks (5, 12 or 20) & No change in outcome \\
\bottomrule
\end{tabular}
\end{table}

The tests also exposed several defects in the live service:
\begin{itemize}
  \item the guided retrieval turn never ran in the deployed service because the sufficiency check
was called without a search function; this was fixed;
  \item a rule containing \texttt{Quelle} also matched inside \texttt{Quelles}, classifying a
procedural question as definitional; the French procedural markers were corrected;
  \item identical documents are now rejected by SHA-256 fingerprint rather than filename, after
859 duplicate chunks had been observed;
  \item an eight-entry concept vocabulary and two additive fields now disclose when an answer was
produced after the request was reinterpreted.
\end{itemize}

At the end of this post-internship update, the offline gate was green (463 unit tests and
121 checks), CI was green, and the question sets and production index were unchanged. Remaining
work concerned topic metadata at chunk level, a topic-retrieval arm, vocabulary facets and a final
measurement campaign.

% ====================== APPENDICES ======================
\appendix

\chapter{Glossary}
\begin{table}[htbp]
\centering
\small
\begin{tabular}{p{3.0cm}p{3.0cm}p{8.5cm}}
\toprule
\textbf{Term} & \textbf{Transliteration} & \textbf{Meaning} \\
\midrule
— & \emph{al-murabaha} & Sale at cost plus a known profit margin \\
— & \emph{al-sigha} & Contractual form (offer and acceptance) \\
— & \emph{al-fasl} & Article in the law; chapter in the circular \\
— & \emph{nafidh} & In force \\
— & \emph{kafin} & Sufficient (sufficiency state) \\
— & \emph{ghayr kafin} & Insufficient \\
— & \emph{mutaarid} & Contradictory (two sources disagree) \\
RAG & — & Retrieval-Augmented Generation \\
\texttt{hit@k} & — & Does the expected passage appear among the top $k$ results? \\
RRF & — & Reciprocal Rank Fusion \\
BM25 & — & A standard lexical ranking function \\
Citation validator & — & Deterministic verification of citation membership and numbers \\
Sufficiency & — & Comparison of evidence requirements with retrieved content \\
Feature gate & — & A component switch with a verified fallback behaviour \\
\bottomrule
\end{tabular}
\end{table}

\emph{Note: Arabic terms are transliterated in this file so it can be compiled with
\texttt{pdflatex} without an Arabic font.}

\chapter{Result Provenance}
\begin{table}[htbp]
\centering
\small
\begin{tabular}{p{5.6cm}p{9.0cm}}
\toprule
\textbf{Result or Data} & \textbf{Source Listed in the Working Repository} \\
\midrule
4 documents, 339 chunks and repairs 215/35/23/21 & \texttt{audits/MASTER\_INDEX.md},
\texttt{AGENTS.md} \\
50 questions (10/20/15/5) and 10 target cases & \texttt{audits/QUESTIONS\_MATRIX.md},
\texttt{audits/PHASE3\_INTERVENTIONS.md} \\
BM25 42/60/62 to 62/82/87; live 51/64/71 to 71/80/87 &
\texttt{audits/PHASE2\_BASELINES.md}, runs 3.1 and 3.2 \\
Vector 71/80/87, RRF 71/80/82 and blend 62/80/84 &
\texttt{audits/PHASE2\_BASELINES.md}, run 3.2 \\
Status agreement 0.720; Jaccard 1.000 & \texttt{audits/PHASE2\_BASELINES.md}, run 3.3 \\
Unit index 100/100/100 vs. 67/92/92 & \texttt{audits/PHASE4\_KNOWLEDGE.md} \\
198 units, 77 relations, 53 numbers and 6 axes & \texttt{audits/PHASE4\_KNOWLEDGE.md} \\
Target categories 50/70/70 & \texttt{audits/PHASE3\_INTERVENTIONS.md} \\
Phase 9: 5/7, 4/5 to 5/5, 3 answers and 3 refusals & \texttt{audits/PHASE9\_STRUCTURE/};
post-internship results \\
Service version at internship end & \texttt{service.py}, \texttt{CONTRACT.md} \\
\bottomrule
\end{tabular}
\end{table}

\chapter{Reproducing the Tests}
\begin{verbatim}
# Environment
cd <repository>/raglab
python3 -m venv .venv && source .venv/bin/activate
pip install -r requirements.txt

# Offline quality gate (no model calls)
./run_tests.sh --offline

# Inspect and index the real corpus, then ask a question
python main.py inspect --data-dir ../docs
python main.py ingest --reset --data-dir ../docs
python main.py answer "..." --query-lang ar

# HTTP service and integration test console
python -m uvicorn service:app
python local_front.py --base-url http://localhost:8000
\end{verbatim}

\chapter{Condensed Timeline of Milestones}
\emph{The available summary log mainly details milestones from 4 September onward. It does not
break down each day from 3 August to 3 September; the table below is therefore not a complete daily
log. The 4--5 October changes are excluded here and appear in the post-internship note.}

\begin{table}[htbp]
\centering
\small
\begin{tabular}{p{2.6cm}p{12.0cm}}
\toprule
\textbf{Period} & \textbf{Main Activities} \\
\midrule
4--6 September & First complete pipeline; multiple providers; query translation; variant fusion;
vector validation; initial measurement harnesses \\
7--8 September & Structural chunking, offline harness, live-model A/B tests, work contract and
interactive console \\
11--16 September & Extraction of the HTTP service and profiles; integration console;
authentication token; output scrubbing; numeric checks; document API and endpoint freeze \\
23--25 September & Chunk browser; reverse specification; gap analysis; production-image packaging \\
28--30 September & Refocusing on the real corpus; document-by-document audits; linguistic repairs;
approval of corrected texts; 50-question test set \\
1 October & Formal baselines; dense-retrieval decision and lexicon deactivation; knowledge and
understanding layers; version 1.3.0 \\
2 October & Sufficiency fields activated; interrogation; remote diagnostics; deployment; token
incident fixed; version 1.4.0 \\
3 October & End of the internship period used in this report \\
\bottomrule
\end{tabular}
\end{table}

\vfill
\begin{center}
\small\emph{Report prepared by \StudentName{} — internship at \HostOrganization, \InternshipPeriod,
supervised by \SupervisorName.}
\end{center}

\end{document}
